\section{Discussion and Conclusion}
\label{sec:discussion}
\label{sec:conclusion}

We set out to separate what an audit reveals from what a detector recovers and from where its alert threshold falls. At an equal expected budget, a probability audit guided by scan rank and label uncertainty finds far more unmarked nodules per review than mark-stratified review, mainly through image ranking and image-dependent uncertainty; the reader marks add a smaller increment. These labels did not raise omission recall: the detector trained on them is a linear head on a frozen embedding, and at the locked threshold almost no omission is detected by any training method (Sec.~\ref{sec:misses}); better labels thus shift the scores of omitted nodules, but rarely beyond the threshold. The audit remains useful in its own right: because every candidate is audited with a known probability, the same sample gives an unbiased estimate of the proportion of positives among the candidates that the first reader left unmarked, whatever detector is trained afterwards, so the completeness of the weak labels can be reported before detectors are compared.

The threshold itself proved to be the more fragile part of the evaluation. With all scores fixed, it moved by several ranks when the \NExcluded{} consensus-positive patients among the \PCal{} calibration negatives were removed, and the estimated gap between the methods moved from $\PairRelativeP$ to $+\PairRelativeU$ points, a shift whose interval excludes zero although the new gap itself does not. At a \AlphaPct\% alert rate with \PCal{} calibration negatives, the threshold is the sixth-largest negative maximum, set by a handful of patients. We recommend that low-alert evaluations of lesion detectors report the reference status of the maxima near the threshold, replay the scores under an alternative negative pool, and compare target recovery and alerts on a common test support (Sec.~\ref{sec:tail}); the negative reference is a modelling decision whose effect should be reported.

Several limitations qualify these conclusions. We measure audit cost in simulated candidate decisions and do not model reading time, additional reads or candidate selection. The roster inherits the LUNA16 proposal sources, some developed on the same LIDC scans. All references reuse the same unblinded LIDC reader records, and no independent adjudication establishes whether the disputed maxima are in-scope nodules. Because LUNA16 is public and extensively studied, the evaluation is retrospective, although all weights, endpoints and contrasts were fixed before the calibration and test subsets were scored; the reference replay and the acquisition ablation were specified after the locked comparison and are follow-up analyses. An end-to-end learner on the audited labels, a sensitivity analysis of Eqs.~\eqref{eq:leverage}--\eqref{eq:allocation}, a sweep over $\alpha$ and replication on independently adjudicated patients are the decisive next steps. This secondary analysis uses released, de-identified data and involves no participant interaction or new data collection.
